\documentclass[10pt,a4paper]{article}
\usepackage{fontspec,amsmath,amssymb,graphicx}
\usepackage[margin=18mm]{geometry}
\setmainfont{Times New Roman}
\setlength{\parindent}{0pt}
\setlength{\parskip}{5pt}
\allowdisplaybreaks
\sloppy
\emergencystretch=2em
\begin{document}
\section*{BT\_5\_1\_3\_08 — Lời giải bổ sung}
(i) Theo $y$, dùng $Y_j=\cos k_jy-\sin k_jy/k_j$, vì $Y'(0)=-Y(0),Y'(1)=0$. Các mode gồm cả khả năng trị riêng âm: viết chung $-Y_j''=\lambda_jY_j$, $d_j=\int_0^1Y_j/\int_0^1Y_j^2$. Mỗi hệ số $V_j$ giải $V_j''-\lambda_jV_j=0$, $V_j(0)=-d_j,V_j(1)=d_j$. Với $\lambda_j>0$,
\[V_j=d_j\frac{\sinh(\sqrt{\lambda_j}x)-\sinh(\sqrt{\lambda_j}(1-x))}{\sinh\sqrt{\lambda_j}}.\]
Với $\lambda_j=-\kappa_j^2<0$, thay sinh bởi sin. Ở đây trị riêng âm duy nhất thỏa $\kappa\tanh\kappa=1$ và $\kappa<\pi$, nên không có cộng hưởng với hai biên Dirichlet theo $x$. $u=\sum V_jY_j$.

(ii) $u=-\sum_{n\text{ lẻ}}4\sin nx\sinh(n(\pi-y))/[\pi n\sinh n\pi]$.

(iii) Miền ghi $(0,\pi)\times(0,1)$ nhưng biên ghi $x=1,y=\pi$: không khép miền, không xác định duy nhất như bản in. Nếu miền dự kiến là $(0,1)\times(0,\pi)$, dùng công thức bốn cạnh Dirichlet với cạnh dưới và phải bằng 1, hai cạnh còn lại 0.

(iv) Đặt $X=x-\pi$. Nghiệm là $1-2v$, với
\[v=\sum_{n\text{ lẻ}}\frac4{\pi n}\sin nX\frac{\sinh ny+\sinh(n(\pi-y))}{\sinh n\pi}.\]

(v) $u=1+w$, $w=0$ ở $x=\pi,2\pi$, $w_y(x,0)=-1,w_y(x,\pi)=0$. Do đó
\[w=\sum_{n\text{ lẻ}}\frac4{\pi n^2}\frac{\cosh(n(\pi-y))}{\sinh n\pi}\sin(n(x-\pi)).\]\end{document}